\begin{recipe}{Verse pasta}
\label{recipe:verse-pasta}
\kicker[Hoofdgerecht,Italiaans,Basisrecept]{Hoofdgerecht · Italiaans · basisrecept}
\index[register]{Verse pasta}
\index[register]{Bloem}
\index[register]{Eidooiers}

\meta{45 minuten \dvd Pastamachine}

\lettrine{M}{et} maar twee ingrediënten maak je een soepel, goudgeel pastadeeg
dat echt anders smaakt dan pasta uit de supermarkt: zachter en rijker door de
dooiers. Niet elke saus is daar even blij mee, verse pasta doet het vooral
goed met een simpele boter- of roomsaus, terwijl een stevige, klonterige saus
beter bij gedroogde pasta past.

\blockrule{Ingrediënten}
\begin{ingredients}
  \ing{4}{eidooiers}\index[register]{Eidooiers}
  \ing{200 g}{bloem type 00 (tipo 00)}\index[register]{Bloem}
  \ingb{extra bloem of stuifbloem, om te bestuiven}
\end{ingredients}

\blockrule{Bereiding}
\tip{Geen pastamachine in huis? Rol het deeg dan met een deegroller flinterdun
uit. Het kost meer moeite, maar het lukt prima.}
\begin{steps}
  \item Doe de bloem op een schone werkbank en maak in het midden een
        kuiltje. Doe de eidooiers erin.
  \item Klop de dooiers met een vork los en meng ze geleidelijk door de bloem
        tot er een grof deeg ontstaat.
  \item Kneed het deeg 8–10 minuten stevig door tot het glad en elastisch
        aanvoelt. Druk zachtjes met een vinger in het deeg: veert het
        kuiltje terug, dan is het elastisch genoeg.
  \item Wikkel het deeg in vershoudfolie en laat 20 minuten rusten op
        kamertemperatuur, zodat de gluten ontspannen en het deeg zich
        makkelijker laat uitrollen.
  \item Verdeel het deeg in stukken en rol elk stuk met de pastamachine steeds
        dunner uit, met tussendoor een beetje stuifbloem zodat het niet
        plakt.
  \item Snijd de uitgerolde lappen tot de gewenste pastavorm.
  \item Kook de pasta 3 minuten in ruim gezouten water, verse pasta is veel
        sneller gaar dan gedroogde.
  \item Giet direct af en meng de pasta meteen door de saus naar keuze.
\end{steps}
\end{recipe}
